\section{Related work}\label{sec:related}

\paragraph{Hawkes processes and their estimation.} Hawkes processes \citep{hawkes1971spectra} admit
a cluster (branching) representation \citep{hawkes1974cluster}, and their likelihood is that of a
general point process with conditional intensity \citep{daley2003introduction}. Exponential and
sum-of-exponential kernels are usually fitted by maximum likelihood \citep{bacry2015hawkes}, for
instance with the \texttt{tick} library \citep{bacry2018tick}. Other estimators exploit the
branching structure through EM \citep{veen2008estimation}, add low-rank and sparse penalties
\citep{zhou2013learning} or group-sparse Granger-causal penalties \citep{xu2016learning}, invert
the second-order statistics through a Wiener--Hopf system
\citep{bacry2012nonparametric,bacry2016first,bacry2016estimation}, match integrated cumulants
\citep{achab2017uncovering} or discretise the process into an INAR model
\citep{kirchner2017estimation}. The branching matrix also defines a graph of causal influence
\citep{embrechts2018hawkes}. None of these estimators returns a certificate of global optimality for
multiscale kernels with delayed peaks. Erlang kernels make a Hawkes process Markovian, which
\citet{duarte2019stability} use to study the stability and simulation of nonlinear Hawkes
processes; we use the same embedding for estimation, certification and closed-form responses.

\paragraph{Hawkes processes in finance.} Branching ratios close to one have been read as evidence of
strong endogeneity \citep{filimonov2012quantifying,hardiman2013critical}, and such nearly unstable
processes have their own limit theory \citep{jaisson2015limit}, although a mis-specified
intraday seasonality biases these estimates upward \citep{filimonov2015apparent,omi2017hawkes}.
Mutually exciting processes also describe contagion across markets \citep{aitsahalia2015modeling}
and lead--lag relations between assets \citep{huth2014high}. \citet{rambaldi2015modeling} add an
exogenous kernel for macroeconomic announcements to a Hawkes model of FX quotes, and
\citet{andersen2003micro} document how announcement surprises move prices. In this literature the
speed of absorption is read off the exogenous kernel. We compute the full response, echoes
included, in closed form, and test the news kernels against matched placebo windows within the
likelihood.

\paragraph{Neural temporal point processes.} \citet{shchur2021neural} and \citet{zhou2026advances}
review the field. RMTPP \citep{du2016recurrent} uses an exponential intensity between events, whose
integral is closed form. NHP \citep{mei2017neural}, SAHP \citep{zhang2020self}, THP
\citep{zuo2020transformer}, AttNHP \citep{yang2022transformer} and the state-space model S2P2
\citep{chang2025deep} compute the intensity with recurrent, attention or state-space encoders, and
estimate the compensator by Monte Carlo or quadrature; continuous-time models driven by neural
differential equations \citep{jia2019neural,chen2021neural} integrate the intensity numerically.
Simulation and prediction usually rely on thinning \citep{ogata1981lewis}. FullyNN
\citep{omi2019fully} models the cumulative hazard with a monotone network, and IntensityFree
\citep{shchur2020intensity} models the density of the next gap with a mixture of log-normals. DTPP
\citep{panos2024decomposable} pairs such a time model with an attention model for the marks. These
models have exact likelihoods but no Hawkes structure. Recent work also explores Mamba-style
encoders \citep{gao2024mamba}. Hyper Hawkes processes \citep{boyd2025hyper} let a network set the
dynamics of a latent linear Hawkes process at each event; a softplus read-out keeps the intensity
positive, so the compensator is again estimated by sampling. EasyTPP \citep{xue2024easytpp}
provides shared implementations, \citet{bosser2023predictive} study next-event prediction, and
HoTPP \citep{karpukhin2024hotpp} evaluates long-horizon forecasts.

EPT-X differs from these models in three ways. First, between events its intensity is a linear
read-out of a phase-type system that can be solved exactly, plus bounded hazard terms and a
renewal channel whose survival function is a mixture of Erlang survival functions. All of these
integrate in closed form, so the likelihood we train on and report is exact. Second, its backbone
is the linear Hawkes process itself, so a fitted model still has excitation kernels that can be
read off, and it can start from a certified MSX fit. Third, its states follow an affine recursion,
so it trains with a parallel scan, as linear state-space layers do \citep{smith2023simplified},
without sampling the compensator.

\paragraph{Phase-type distributions.} Phase-type distributions are the absorption times of finite
Markov chains, and they can be fitted by EM \citep{asmussen1996fitting}. Mixtures of Erlang
distributions are dense in the set of distributions on $(0,\infty)$
\citep{tijms1994stochastic,lee2010modeling}. We use this fact for the renewal channel, and the
Erlang-chain recursion for every state in the paper.
